%%%%%%%%%%%%%%%%%%%%%%%%%%%%%%%%%%%%%%%%%%%%%%%%%%%%%%%%%%%%%%%%%%%%%%
% Lớp 12 | Tọa độ của véc-tơ trong không gian | Tọa độ của véc-tơ và các công thức
% Dạng 2H2?2-1: Công thức lý thuyết
% Số câu: 4  (sinh tự động bởi tach_cau.py; từ nay thầy sửa trực tiếp file này)
%%%%%%%%%%%%%%%%%%%%%%%%%%%%%%%%%%%%%%%%%%%%%%%%%%%%%%%%%%%%%%%%%%%%%%

%%%%%%%%%%%%%%%% TRẮC NGHIỆM 4 PHƯƠNG ÁN %%%%%%%%%%%%%%%%

% nguồn: Lop_12_Hinh_hoc_Chuong_2
\begin{ex}%[2H2N2-1]%[ma:3867159b]
	Trong không gian $Oxyz$,  tọa~độ của véc-tơ $\vec{u}=x \vec{i}+y \vec{j}+z \vec{k}$~là
	 \choice
		{$(y;  z; x)$}
		{$(z; x;  y)$}
		{\True $(x; y; z)$} 
		{$(x;  z; y)$}
	\loigiai{ 
		Ta có  $\vec{u}=(x;y;z)$.
	 }
\end{ex}

% nguồn: Lop_12_Hinh_hoc_Chuong_2
\begin{ex}%[2H2N2-1]%[ma:653295d4]
	Trong không gian $Oxyz$, cho $\vec{u}=(x;  y; z)$ khác véc-tơ-không.\xuongdong
	Với $k \in \mathbb{R} \setminus\{0\}$, tọa~độ  của véc-tơ $k \vec{u}$~là
	\choice 
		{\True $(kx; ky; kz)$} 
		{$(kz;  kx; ky)$}
		{$(ky; kz; kx)$} 
		{$(kx;  kz; ky)$}
	 \loigiai{
		Ta có $k\vec{u}=(kx;ky;kz)$, với $k \in \mathbb{R} \setminus\{0\}$. 
	 }
\end{ex}

% nguồn: Lop_12_Hinh_hoc_Chuong_2
\begin{ex}%[2H2N2-1]%[ma:50fb6566]
	Trong không gian $Oxyz$, cho hai véc-tơ $\vec{u}=(x; y; z)$, $\overrightarrow{u'}=\left(x';  y'; z'\right)$. Tọa độ của véc-tơ $\vec{u}+\overrightarrow{u'}$~là
	\choice 
		 {\True $\left(x+x'; y+y'; z+z'\right)$}
		{$\left(x-x'; y-y';  z-z'\right)$}
		{$\left(x'-x;  y'-y; z'-z\right)$}
		{$\left(x  x'; yy'; zz'\right)$}
	 \loigiai{
		Ta có $\vec{u}+\overrightarrow{u'}=\left(x+x'; y+y';  z+z'\right)$.
	} 
\end{ex}

% nguồn: Lop_12_Hinh_hoc_Chuong_2
\begin{ex}%[2H2H2-1]%[ma:8378c527]
	Trong không gian $Oxyz$,  cho điểm $M(1;2;3)$. Gọi $M_1$, $M_2$ lần lượt là hình chiếu vuông góc của $M$ lên các trục $Ox$, $Oy$. Véc-tơ nào dưới đây là một véc-tơ chỉ phương của đường thẳng $M_1M_2$?
	\choice 
		{\True  $\overrightarrow{u}_4=(-1;2;0)$}
		 {$\overrightarrow{u}_1=(0;2;0)$}
		{$\overrightarrow{u}_2=(1;2;0)$} 
		 {$\overrightarrow{u}_3=(1;0;0)$}
	\loigiai{Hình chiếu vuông góc của $M$ lên trục $Ox$ là $M_1(1;0;0)$.\\ 
		Hình chiếu vuông góc của $M$  lên trục $Oy$ là $M_2(0;2;0)$.\\
	Véc-tơ  chỉ phương của đường thẳng $M_1M_2$ là $\overrightarrow{M_1M_2}=(-1;2;0)$.}
\end{ex}
